\chapter{The Thesis}
\label{ch:thesis}

\section{Three wars, one playbook}

Between roughly 2005 and 2025, China fought and won three complete technology wars. In solar photovoltaics, it went from assembler of imported cells to producing over 90\% of every stage of the global supply chain, while module prices fell by more than 99\% from their 1976 level~\cite{owid-learning,pv-shares-2023}. In lithium-ion batteries, a technology commercialized by Sony in 1991 and led for two decades by Japan and Korea, Chinese firms now hold roughly 69--75\% of global EV battery installations and over 80\% of global cell manufacturing capacity, with pack prices down 93\% in real terms since 2010~\cite{sne-2025,iea-gevo-2026,bnef-2025}. In electric vehicles, China became simultaneously the world's largest car market, the world's largest car exporter, and home to the world's largest EV maker, while foreign incumbents' share of the Chinese market collapsed from 62\% to 35\% in five years~\cite{rhodium-evs,caam-2025}.

These were not three lucky outcomes. They were three executions of a single, learnable playbook, whose elements Chapter~\ref{ch:playbook} distills: patient state designation of a strategic industry; demand creation at home; subsidized scale-up of manufacturing far ahead of demand; brutal domestic competition that functions as an evolutionary selection engine; capture of the entire upstream supply chain; a cost-down learning-curve war that bankrupts foreign incumbents; and finally, once dominance is achieved, the conversion of that dominance into geopolitical leverage---export controls on the very technologies China was once accused of copying~\cite{mofcom-battery-controls,csis-rare-earth}.

\section{Why AI is the fourth war}

The prevailing Western assumption is that AI is different: that frontier AI depends on scarce, controllable inputs---leading-edge logic, high-bandwidth memory (HBM: stacked DRAM mounted on the accelerator package), proprietary interconnects, closed model weights---and that export controls on these chokepoints can durably preserve a lead. This book argues the assumption is failing on all three fronts simultaneously, and failing in a characteristic way: not because China matches the West chokepoint-for-chokepoint, but because it is engineering the chokepoints out of the system.

\begin{description}
  \item[The model (Chapter~\ref{ch:model}).] Chinese laboratories have converted the model layer into a commons. As of mid-2026, Chinese open-weight models occupy the top open positions on independent intelligence indices, account for roughly 41\% of all Hugging Face downloads and a majority of tokens routed through neutral API aggregators (application programming interfaces---metered, hosted access to a model), and are priced at a small fraction of Western closed models---while the measured capability gap to the best closed US models has compressed to low single digits~\cite{hf-spring-2026,openrouter-2026,stanford-aiindex-2026}. Open weights at marginal cost do to the model layer what Chinese modules did to solar: they set a world price that rent-seeking business models cannot survive.
  \item[The compute (Chapter~\ref{ch:compute}).] The state has designated the open RISC-V instruction set architecture (ISA---the contract between software and processor, and the layer at which Arm and x86 collect licence rents) as national infrastructure, is standardizing matrix/vector extensions and chip-to-chip interconnects across vendors, and has open-sourced the software stacks (CANN, MindSpore, UnifiedBus specifications) that would otherwise fragment the domestic ecosystem~\cite{riscv-policy,huawei-cann,huawei-ub}. In parallel, a memory-capacity arbitrage---substituting terabytes of cheap LPDDR (low-power DDR, the commodity mobile-class DRAM of phones and laptops) for scarce HBM---attacks the assumption that frontier training requires centralized, HBM-bound superclusters at all (Section~\ref{sec:lpddr6}).
  \item[The semiconductor (Chapter~\ref{ch:semi}).] Beneath both layers, the fabs. SMIC has reached 5\,nm-class logic without EUV; CXMT has become the world's fourth DRAM producer and the swing supplier in a global memory shortage; the domestic equipment industry has placed three firms in the global top twenty~\cite{smic-5nm,cxmt-ipo,toolmakers-top20}. The remaining gaps---EUV lithography, leading-edge HBM---are real, but the strategy does not require closing them on the West's terms; it requires making them irrelevant to the chosen architecture.
\end{description}

\section{Openness as a weapon}

The connective tissue among the three elements is a deliberate inversion of the Western technology-rent model. Where the United States concentrated AI value in proprietary layers---CUDA, NVLink, closed frontier weights, x86/Arm licensing---China's strategy, formalized in the July 2025 Global AI Governance Action Plan and the ``AI+'' State Council initiative, is to make every one of those layers open, standardized, and free~\cite{action-plan,ai-plus}. This is not altruism; it is the same move as pricing solar modules at \$0.09 per watt. An open layer generates no rent for the incumbent, but it generates \emph{adoption}, and adoption generates the scale, the ecosystem, and the standards-setting power that constitute victory in the precedent industries. The West's chokepoint strategy and China's commoditization strategy are not symmetric opponents: one defends margins, the other destroys them.

\section{Four propositions, stated separately}
\label{sec:propositions}

Rigor requires decomposing the thesis into four propositions that are related but logically independent---none proves the next, and each has its own evidence and its own failure modes:

\begin{description}
  \item[P1 --- Adoption.] China is winning open-weight model adoption. \emph{Confirmed by}: download, derivative, and routed-token majorities persisting and extending into enterprise spend. \emph{Falsified by}: a durable reversal of the OpenRouter/Hugging Face shares, or enterprise open-weight spend stagnating below $\sim$15\% while usage metrics plateau. Current grade: strongly supported [V].
  \item[P2 --- Sovereign compute.] China is constructing a domestically sufficient training-and-serving stack. \emph{Confirmed by}: a frontier-class model trained end-to-end on domestic silicon, packaging, and memory. \emph{Falsified by}: persistent dependence on stockpiled or smuggled foreign accelerators and HBM past $\sim$2028. Current grade: rapid progress, not yet achieved [P/A/M].
  \item[P3 --- Rent destruction.] Open-weight commoditization will compress the closed-frontier rents that finance Western laboratories. \emph{Confirmed by}: closed-model gross revenue growth decelerating while capability parity holds; margin compression propagating up the stack. \emph{Falsified by}: closed labs sustaining pricing power through service, distribution, and trust layers even at model parity. Current grade: mechanism operating, magnitude contested [V/S].
  \item[P4 --- Financial repricing.] The compression of P3, if it occurs, punctures the US AI capital structure. \emph{Confirmed by}: the indicator table of Chapter~\ref{ch:trade}. \emph{Falsified by}: demand elasticity (the Jevons case, \S\ref{sec:jevons}) absorbing the price decline---AI getting cheap faster than usage grows is the bear case; usage growing faster than prices fall is the bull case, and it rescues the buildout even under P1--P3. Current grade: scenario judgment [S].
\end{description}

Figure~\ref{fig:props} shows the propositions with their couplings and their contrary evidence. P1 can hold while P2 fails (China dominating distribution on stockpiled NVIDIA silicon); P3 can hold while P4 fails (margins compress while volume growth fills the data centers); and a US valuation correction could arrive from macro causes with none of P1--P3 mattering. The convergence argument of this book is that P1 is established, P2 is trending, and the \emph{coupling} of P1$\times$P2---open weights meeting sovereign, capacity-first hardware---is what forces P3, with P4 as its priced-in consequence. Chapters~\ref{ch:model}--\ref{ch:trade} take the propositions in order.

\begin{figure}[htbp]
\centering
\begin{tikzpicture}[
  every node/.style={font=\scriptsize},
  prop/.style={draw=inkMut, line width=0.7pt, rounded corners=3pt, align=left,
               inner sep=5pt, text width=30mm, minimum height=17mm},
  cpl/.style={-{Stealth[length=4pt]}, line width=1pt, color=inkSec},
  ev/.style={font=\scriptsize\color{inkSec}, align=center, text width=30mm},
]
\node[prop, fill=sOne!10, draw=sOne] (p1) at (0,0)
  {\textbf{P1 Adoption}\\open-weight distribution and usage\\[2pt]\textbf{[V]} supported};
\node[prop, fill=sThree!10, draw=sThree, right=6mm of p1] (p2)
  {\textbf{P2 Sovereign compute}\\domestically sufficient stack\\[2pt]\textbf{[P/A/M]} trending};
\node[prop, fill=sTwo!10, draw=sTwo, right=6mm of p2] (p3)
  {\textbf{P3 Rent destruction}\\closed-frontier margins compress\\[2pt]\textbf{[V/S]} operating};
\node[prop, fill=sFour!10, draw=sFour, right=6mm of p3] (p4)
  {\textbf{P4 Repricing}\\the capital structure breaks\\[2pt]\textbf{[S]} scenario};

\draw[cpl] (p1) -- (p2);
\draw[cpl] (p2) -- (p3);
\draw[cpl] (p3) -- (p4);
\draw[cpl, dashed] (p1.north) to[out=40,in=140] node[above, font=\scriptsize\color{inkSec}] {price floor} (p3.north);

\node[ev, below=3mm of p1] {evidence: downloads, derivatives, routed tokens; \emph{not} enterprise spend};
\node[ev, below=3mm of p2] {evidence: LineShine, Ascend/CXMT ramps; \emph{not} a domestic frontier training run};
\node[ev, below=3mm of p3] {evidence: token deflation, migrations; \emph{contra}: service-layer rents, elasticity};
\node[ev, below=3mm of p4] {evidence: leverage, concentration; \emph{contra}: real revenue growth};

\node[font=\scriptsize\color{inkSec}, align=left, text width=146mm, below=25mm of p2.south, xshift=13mm]
  {\textbf{None of these arrows is entailment.} P1 can hold while P2 fails (dominance on stockpiled foreign
   silicon); P3 can hold while P4 fails (margins compress while volume growth fills the data centres, the
   elasticity case of \S\ref{sec:jevons}); and P4 could arrive from macro causes with none of P1--P3 mattering.
   The convergence argument of this book is that the \emph{coupling} P1$\times$P2---open weights meeting sovereign,
   capacity-first hardware---is what forces P3, with P4 as its priced-in consequence.};
\end{tikzpicture}
\caption[The four propositions and their couplings]{The argument, disaggregated. Each proposition carries its own
evidence grade, its own confirming metrics, and its own contrary evidence; the book's thesis is a claim about the
coupling between them rather than a single chain of entailment. Grades as of the version stamp; Appendix~\ref{app:evidence}
carries the full ledger and Appendix~\ref{app:dashboard} the quarterly indicators that move them.}
\label{fig:props}
\end{figure}


\section{Evidence grades}
\label{sec:evidence}

A claim's evidentiary weight is never self-evident from its phrasing, and so load-bearing claims are graded throughout: \textbf{[V]} independently verified (teardown, benchmark submission, audited filing, multiple independent confirmations); \textbf{[P]} primary-source claim (a company or government statement, reliable as to what was claimed, not what was achieved); \textbf{[A]} analyst estimate; \textbf{[M]} author-modeled result (including companion feasibility studies); \textbf{[R]} roadmap; \textbf{[S]} scenario judgment. Appendix~\ref{app:evidence} collects the headline claims in one graded ledger. Where a section rests on the author's own companion studies---notably the memory-architecture analyses of Chapter~\ref{ch:compute}---the text says so explicitly: those are modeled architectures, not measured production systems.

The thesis, stated in its conditional form: the United States retains the absolute model frontier, the leading semiconductor and cloud ecosystem, and vastly greater private capital. China increasingly leads a different contest---the commoditization of frontier-adjacent intelligence and the construction of a sovereign, though incomplete, compute stack. The decisive question is whether Chinese cost reduction, adoption, standards formation, and hardware sufficiency compound faster than US capability, distribution, and monetization---and whether China closes its remaining gaps in memory, packaging, tools, software maturity, and trust before the open-model economy compresses the rents that finance the Western frontier. This book argues that convergence is the way to bet; the chapters that follow document why, and Chapter~\ref{ch:synthesis} states what would change the bet.

The reader who finishes Part~I will recognize, in Part~II, every move already made once, twice, or three times before---and, in Chapter~\ref{ch:analogy}, the places where the precedents genuinely do not transfer.
